\documentclass[../../DocumentoCompleto.tex]{subfiles}
\begin{document}
\chapter{Definición del concepto}
\section{Introducción}
HOLA
Para el desarrollo de un producto que logre satisfacer las necesidades del cliente, resulta indispensable implementar en primera instancia la fase de generación de conceptos. Donde es una etapa se fundamenta en representaciones preliminares que integran un bosquejo aproximado acompañado de una breve descripción funcional, con el fin de visualizar de manera temprana cómo podría materializarse y funcionar la solución propuesta \parencite[120]{Ulrich2020}.

La relevancia de esta fase se basa en anticipar de qué manera el diseño responderá a los requerimientos establecidos, al tiempo que proporciona una base sólida para orientar las decisiones posteriores. La calidad del concepto inicial es determinante, ya que de ella depende que el producto final cumpla con las expectativas del usuario y alcance un desempeño satisfactorio rigiéndose a las necesidades previamente establecidas \parencite[120]{Ulrich2020}.

En consecuencia, se hace necesario realizar una búsqueda amplia y rigurosa de alternativas, analizando la utilidad, limitaciones y potencial de cada una. Restringirse a una o dos opciones reduce las posibilidades de innovación y aumenta el riesgo de cometer errores en etapas posteriores; por el contrario, la evaluación de múltiples alternativas favorece la elección de la propuesta más adecuada, disminuye la probabilidad de retrasos y reduce de manera significativa los riesgos de fracaso durante la implementación \parencite[120]{Ulrich2020}.


\section{Metodología}

Para llevar a cabo una adecuada generación de conceptos es necesario considerar un proceso estructurado compuesto por cinco fases: (1) aclarar el problema, que permite definir el alcance y comprender las necesidades del usuario; (2) buscar externamente, orientado a identificar soluciones y referentes en el entorno;  (3) buscar internamente, donde se fomenta la creatividad y aportes del equipo de diseño; (4) explorar sistemáticamente, que organiza y combina las alternativas generadas; y finalmente, (5) reflexionar sobre las soluciones y el proceso, con el fin de evaluar los resultados obtenidos y orientar las decisiones posteriores \parencite[126]{Ulrich2020}.


\begin{figure}[h!]
    \centering
    \includegraphics[width=0.9\textwidth]{Figuras/GENERACIONCONCEPTOS.png}
    \caption{\textit{Proceso de generación de conceptos.}}
    \label{fig:procesodegeneraciondeconceptos}
\end{figure}



  \subsection{Aclarar el problema}
  El primer paso en la generación de conceptos es aclarar el problema, lo que consiste en desarrollar un entendimiento general y, si es necesario, descomponerlo en subproblemas. Las entradas ideales para este proceso son la declaración de la misión del proyecto, la lista de necesidades del cliente y las especificaciones preliminares del producto. En muchos casos, el problema se puede dividir en subproblemas más sencillos mediante una descomposición funcional, representando primero el producto, el cual es definido como una “caja negra”, donde es definida como una representación del producto que muestra únicamente sus entradas y salidas de material, energía señales, o parámetros esenciales, sin detallar cómo se realiza la función interna, permitiendo así definir el problema de manera general antes de dividirlo en subfunciones \parencite[122]{Ulrich2020}.


\subsection{Buscar externamente}
El segundo paso en la generación de conceptos es buscar externamente, lo cual implica identificar soluciones ya existentes en el campo de la salud, tanto para el problema general como para los subproblemas definidos. Este proceso se fundamenta en recopilar información de diversas fuentes, con el fin de aprovechar conocimientos previos y así reducir reinvenciones innecesarias \parencite[126]{Ulrich2020}.

\subsection{Buscar internamente}
Posteriormente, se lleva a cabo la búsqueda interna, la cual consiste en generar nuevas ideas de solución al problema con el fin de ampliar las alternativas y enfocarse en producir la mayor cantidad posible de propuestas, incluso aquellas que puedan parecer poco factibles. En este proceso resulta útil el uso de recursos gráficos, como bosquejos, imágenes o maquetas, que permiten explorar diversas posibilidades para ser posteriormente evaluadas y refinadas. En esta etapa se definieron los conceptos de material, aditivo y geometría \parencite[130]{Ulrich2020}.

\subsection{Explorar sistemáticamente}
En la fase de exploración sistemática, se debe organizar y sintetizar toda la información recopilada en las búsquedas internas y externas. Para ello, se utiliza un árbol de clasificación de conceptos, que permite identificar y comparar las opciones más viables y descartar aquellas menos prometedoras. De igual manera, se recurre a una tabla de combinación de conceptos, la cual facilita la selección y evaluación de alternativas de manera ordenada. Cabe recalcar que esta etapa no busca definir la solución final, sino organizar la información obtenida previamente \parencite[133]{Ulrich2020}.


\subsection{Reflexionar sobre las soluciones y el proceso}
La etapa final del proceso de generación de conceptos se basa en realizar un análisis crítico del trabajo desarrollado, con el propósito de identificar aciertos, errores y oportunidades de mejora. Esto permite evaluar si la descomposición del problema, así como las búsquedas externas e internas, se llevaron a cabo correctamente. De igual manera, en esta fase se reconocen posibles sesgos en las ideas, con el fin de fortalecer la calidad del resultado final. \parencite[133]{Ulrich2020}.


 \subsection{Filtrado de los conceptos}
  Una vez organizadas las alternativas, se evaluó la solución más viable empleando una matriz de filtrado y una matriz de evaluación de conceptos.
  
   \subsubsection{Matriz de filtrado}
La matriz de filtrado es una herramienta que sirve para comparar alternativas según criterios definidos y descartar rápidamente las menos viables, conservando solo las opciones con mayor potencial, donde a cada criterio se le asignó una valoración de 1 cuando cumplía de manera positiva, 0 si no tenía efecto significativo y –1 cuando presentaba un aspecto negativo o perjudicial; Posteriormente, se sumaron los resultados para obtener una evaluación neta que permitió identificar cuáles conceptos cumplían de forma más adecuada con los requisitos definidos.


    \subsubsection{Matriz de evaluación de concepto}
La matriz de evaluación es una herramienta utilizada después del filtrado inicial para comparar de manera más detallada los conceptos que continúan en análisis. A diferencia de la matriz de filtrado, esta no solo descarta, sino que asigna calificaciones ponderadas a cada alternativa según criterios previamente definidos (como citotoxicidad, proliferación celular, cicatrización, etc.), considerando además el peso relativo de cada criterio; su finalidad es jerarquizar de forma cuantitativa y objetiva las opciones restantes, permitiendo identificar la propuesta que ofrece el mejor equilibrio entre desempeño, viabilidad y alineación con los objetivos del proyecto.

 
  \subsection{Recomendación del concepto}
Se seleccionó una propuesta para el material, aditivo y geometría respectivamente, teniendo en cuenta una reflexión sobre el concepto que obtuvo la mayor puntuación en cada matriz de evaluación, en la cual se analizaron su alcance y los riesgos asociados a su implementación, lo que permitió confirmar la solidez de la alternativa elegida.


\section{Resultados}



  \subsection{Concepto de material}


    \subsubsection{Hidrogel gelatina-colágeno}
Combina gelatina derivada del colágeno y colágeno tipo I. Se fabrica mediante impresión 3D, brindando una porosidad y arquitectura personalizable para favorecer la retención de humedad y la adaptabilidad anatómica. Ofrece buen soporte celular, estimula la regeneración tisular y mantiene un ambiente húmedo ideal. \cite{Li2022}



\begin{figure}[h!]
	\centering
	\includegraphics[width=0.4\textwidth]{Figuras/HIDROGELGELATINACOLAGENO.png}
	\caption{\textit{Representación de hidrogel impreso en 3D.} 
		Tomado de \textit{“3D-Printed Gelatin-Alginate Hydrogel Dressings for Burn Wound Healing: A Comprehensive Study”} 
		\parencite{Fayyazbakhsh2022}.}
	\label{fig:hidrogel}
\end{figure}



    \subsubsection{Piel de Tilapia}
Material biológico, rico en colágeno tipo I y estructura similar a la dermis humana. Ha mostrado reducción del dolor, menor frecuencia de curaciones y aceleración de la cicatrización. El material se prepara mediante limpieza, esterilización (química o gamma) y conservación en glicerol. Es económica y especialmente útil en regiones con producción acuícola. \cite{ArauzMadrigal2022}



\begin{figure}[h!]
	\centering
	\includegraphics[width=0.4\textwidth]{Figuras/PIEL DE TILAPIA.png}
	\caption{\textit{Representación del uso de piel de tilapia en el tratamiento de quemaduras.}
		Elaboración propia basada en \textit{“300+ burn victims treated with tilapia skin in Ceará since 2016”} 
		\parencite{Sousa2021}.}
	\label{fig:piel_tilapia}
\end{figure}



    \subsubsection{Colágeno}    
El colágeno de tipo I, es utilizado como matriz estructural por su alta biocompatibilidad, bajo riesgo inmunológico y capacidad para promover granulación tisular y migración celular. Puede derivarse de tejidos animales y ser purificado industrialmente. Tienen mejor desempeño clínico frente a apósitos convencionales, reduciendo el tiempo de curación y la necesidad de injertos. \cite{Singh2011}



    \subsubsection{Quitosano}   
Polímero natural derivado de la quitina, el quitosano posee propiedades antimicrobianas, hemostáticas y bioadhesivas. Es soluble en medios ácidos y puede procesarse en forma de películas, espumas o nanofibras. Favorece la cicatrización mediante la estimulación de células epiteliales y fibroblastos, además de actuar como barrera contra infecciones. \cite{Mohandas2015}



    \subsubsection{Alginato}   
Derivado de algas pardas, el alginato forma hidrogeles al contacto con calcio. Es altamente absorbente, biocompatible y mantiene un entorno húmedo ideal para la cicatrización. Se usa principalmente en apósitos para heridas exudativas, aunque por sí solo no promueve la regeneración tisular activa. Es fácil de manipular y económico. \cite{Boateng2008}

    \subsubsection{Agarosa}
Polisacárido neutro extraído de algas rojas, utilizado frecuentemente como matriz inerte en cultivos celulares y apósitos. Aunque su bioactividad es limitada, ofrece excelente retención de humedad y propiedades mecánicas básicas. Se gelifica fácilmente al enfriar, permitiendo su moldeado en apósitos blandos. \cite{Lee2012}

    \subsubsection{Nanocompuestos de quitosano y plata}
Este sistema combina quitosano (biocompatible y bioadhesivo) con nanopartículas de plata que aportan potente actividad antimicrobiana. Garantiza una liberación controlada de plata, buena estabilidad estructural y ausencia de citotoxicidad significativa. Se fabrica mediante dispersión de las plata en solución de quitosano, seguida de gelificación o moldeado. \cite{Islam2023} 



    \subsubsection{Nanofibras electrohiladas de quitosano-PVP}
Producidas por electrohilado, generando fibras delgadas con gran superficie de contacto. El quitosano aporta biocompatibilidad, el PVP estabiliza la estructura y la florizina que es un antioxidante natural modula la inflamación y el estrés oxidativo. Aceleran la epitelización y controlan la proliferación bacteriana. \cite{Zhang2022} 




\begin{figure}[h!]
	\centering
	\includegraphics[width=0.4\textwidth]{Figuras/NANOFIBRAS ELECTROHILADAS DE QUITOSANO PVP.png}
	\caption{\textit{Representación de nanofibras electrohiladas.}
		Tomado de \textit{“Antibacterial and antioxidant phlorizin-loaded nanofiber film effectively promotes the healing of burn wounds”} 
		\parencite{Yang2024}.}
	\label{fig:nanofibras}
\end{figure}



    \subsubsection{Membrana Amniotica}
Tejido de origen humano con propiedades antiinflamatorias, antibacterianas y regenerativas. Rico en colágeno, factores de crecimiento y citoquinas. 
Se conserva en glicerol al 85\% y se aplica directamente sobre la herida. 
En quemaduras de segundo grado, ha mostrado reducción del dolor, disminución del riesgo de infección y mejor resultado estético final; sin embargo, su obtención es compleja e involucra principios éticos. \cite{Gaviria2018}




\begin{figure}[h!]
	\centering
	\includegraphics[width=0.4\textwidth]{Figuras/MEMBRANA AMNIOTICA.png}
	\caption{\textit{Representación de membrana amniótica utilizada como biomaterial.}
		Tomado de \textit{“Placentas: de la basura al tejido que lo 'repara todo”} 
		\parencite{RamirezDeCastro2024}.}
	\label{fig:membrana_amniotica}
\end{figure}




    \subsubsection{Nanofibras de policaprolactona, gelatina, poliestireno, colágeno y quitosano}
Diseño biomimético multicapa inspirado en la hoja de loto. Presenta una capa externa hidrofóbica (PCL/poliestireno) y una capa interna hidrofílica (gelatina, colágeno, quitosano). Se fabrica mediante electrohilado multicapa. Esta combinación simula las funciones selectivas de la piel: protege contra contaminantes y favorece la regeneración en el interior además de acelerar la curación y tener control bacteriano. \cite{Wang2023}




\begin{figure}[h!]
	\centering
	\includegraphics[width=0.35\textwidth]{Figuras/NANOFIBRAS DE POLICAPROLACTONA GELATINA POLIESTIRENO COLAGENO QUITOSANO.png}
	\caption{\textit{Representación de nanofibras biomiméticas.}
		Tomado de \textit{“Nanofibers: An effective biomedical tool for burn management”} 
		\parencite{Sen2023}.}
	\label{fig:nanofibras_biomimeticas}
\end{figure}




  \subsection{Concepto de aditivo}
  
    \subsubsection{Gentamicina}
Antibiótico aminoglucósido utilizado en quemaduras infectadas, eliminando bacterias gram negativas al inhibir la síntesis de proteínas esenciales para su crecimiento. \cite{MarmolejoRuiz2025}

   \subsubsection{Minociclina}
Antibiótico de amplio espectro útil en quemaduras infectadas, bloquea la síntesis proteica bacteriana y reduce el riesgo de complicaciones infecciosas. \cite{DelgadoBeltran2022}


   \subsubsection{Vancomicina}
Antibiótico glucopéptido indicado en quemaduras complicadas con infecciones por bacterias grampositivas resistentes, impide la formación de su pared celular protectora. \cite{Teixeira2023}


   \subsubsection{Sulfadiazina de plata}
Crema tópica con efecto bactericida y bacteriostático que previene infecciones en quemaduras, liberando plata iónica sobre el área afectada. \cite{CardenasZuluaga2023}


   \subsubsection{Colagenasa}
Enzima desbridante que elimina tejido necrótico en quemaduras, favoreciendo la limpieza de la herida y acelerando la regeneración del tejido. \cite{MonteroUscanga2021}

   \subsubsection{Fitostimoline}
Crema tópica vegetal que favorece la cicatrización en quemaduras, reduce inflamación y acelera la regeneración dérmica gracias a su extracto de trigo. \cite{AjitimbayLopez2023}

   \subsubsection{Ácido hialurónico}
Componente natural que hidrata, protege y estimula la regeneración dérmica en quemaduras, mejorando elasticidad y favoreciendo la cicatrización de heridas. \cite{RituertoCuerdo2023} 

   \subsubsection{Aloe vera}
Gel vegetal con acción antiinflamatoria, calmante y cicatrizante, usado para aliviar dolor, reducir inflamación y regenerar piel quemada. \cite{CabreraVergara2024}

   \subsubsection{Caléndula}
Extracto floral con propiedades antiinflamatorias, antisépticas y cicatrizantes, indicado para tratar quemaduras leves y prevenir infecciones superficiales. \cite{SanchezFlorez2023}


  \subsection{Concepto de geometría}

  \subsubsection{Geometría tradicional de primer orden }
  
Estructura formada por cortes paralelos cortos y estrechos para generar una malla estándar. Al estirarse forma aperturas en forma de rombos. Esta geometría proporciona cobertura con baja capacidad de expansión alcanzando una relación de mallado de 1.5:1. \cite{GuptaChanda2023} 



\begin{figure}[h!]
	\centering
	\includegraphics[width=0.4\textwidth]{Figuras/TRADICIONAL PRIMER ORDEN.png}
	\caption{\textit{Geometría tradicional de primer orden.}
		Tomado de \textit{“Análisis de elementos finitos de patrones jerárquicos basados en metamateriales para generar alta expansión en injertos de piel”} 
		\parencite{GuptaChanda2023}.}
	\label{fig:geometria_primer_orden}
\end{figure}



 \subsubsection{Geometría tradicional de segundo orden}
Geometría formada por el mismo patrón que la estructura de primer orden pero con mayor distancia entre cortes o cortes paralelos más largos. Al estirarse forman apertura en forma de elipse y alcanza una 
relación de mallado de hasta 3:1 proporcionando una expansión moderada. \cite{GuptaChanda2023}



\begin{figure}[h!]
	\centering
	\includegraphics[width=0.4\textwidth]{Figuras/TRADICIONAL SEGUNDO ORDEN.png}
	\caption{\textit{Geometría tradicional de segundo orden.}
		Tomado de \textit{“Análisis de elementos finitos de patrones jerárquicos basados en metamateriales para generar alta expansión en injertos de piel”} 
		\parencite{GuptaChanda2023}.}
	\label{fig:geometria_segundo_orden}
\end{figure}



\subsubsection{Geometría tradicional de tercer orden }
Estructura formada por cortes que forman una malla estándar con incisiones largas y mayores separaciones comparado con las geometrías de orden 1 y 2. Al expandirse genera una red abierta alcanzando en aplicación clínica una relación de mallado de 3.4:1. \cite{GuptaChanda2023} 



\begin{figure}[h!]
	\centering
	\includegraphics[width=0.4\textwidth]{Figuras/TRADICIONAL TERCER ORDEN.png}
	\caption{\textit{Geometría tradicional de tercer orden.}
		Tomado de \textit{“Análisis de elementos finitos de patrones jerárquicos basados en metamateriales para generar alta expansión en injertos de piel”} 
		\parencite{GuptaChanda2023}.}
	\label{fig:geometria_tercer_orden}
\end{figure}


\subsubsection{Geometría auxética tipo “Y”}

Estructura formada por patrones en forma de “Y” conectadas en estrella. Permite una buena expansión bajo esfuerzos moderados, alcanzando una relación de mallado de 3:1 proporcionando cobertura y resistencia.\cite{Gupta2023} 

\begin{figure}[h!]
	\centering
	\includegraphics[width=0.4\textwidth]{Figuras/TIPO Y.png}
	\caption{\textit{Geometría auxética tipo “Y”.}
		Tomado de \textit{“High expansion auxetic skin graft simulants for severe burn injury mitigation”} 
		\parencite{Gupta2023}.}
	\label{fig:geometria_auxetica_y}
\end{figure}


\subsubsection{Geometría auxética tipo “I”}

Geometría formada por patrones en forma de “I” alineadas verticalmente. Esta estructura se caracteriza por una expansión uniforme con bajo esfuerzo y alta resistencia logrando una relación de mallado superior a 4:1 lo que indica el cubrimiento de áreas extensas. \cite{Gupta2023} 


\begin{figure}[h!]
	\centering
	\includegraphics[width=0.4\textwidth]{Figuras/TIPO I.png}
	\caption{\textit{Geometría auxética tipo “I”.}
		Tomado de \textit{“High expansion auxetic skin graft simulants for severe burn injury mitigation”} 
		\parencite{Gupta2023}.}
	\label{fig:geometria_auxetica_i}
\end{figure}


\subsubsection{Geometría auxética tipo panal de abejas}
Geometría basada en un patrón formado por estructuras hexagonales conectadas por secciones delgadas. Tiene una expansión estimada de 267\% cuando hay una distribución de esfuerzo óptima. Sin embargo, la relación de mallado es de aproximadamente 2.2:1 siendo menor comparada con las otras geometrías. \cite{Gupta2023} 


\begin{figure}[h!]
	\centering
	\includegraphics[width=0.4\textwidth]{Figuras/TIPO PANAL DE ABEJAS.png}
	\caption{\textit{Geometría auxética tipo panal de abejas.}
		Tomado de \textit{“High expansion auxetic skin graft simulants for severe burn injury mitigation”} 
		\parencite{Gupta2023}.}
	\label{fig:geometria_auxetica_panal}
\end{figure}


\subsubsection{Geometría auxética tipo “H”}
Geometría con patrones en forma de “H” las cuales al aplicar tensión se abren en cuatro vacíos cuadrados. Esta estructura tiene una mayor expansión de área vacía alcanzado hasta el 1032\% pero una baja relación de mallado de 1.8:1 por lo cual no es ideal para garantizar cobertura, además puede presentar riesgo de fractura. \cite{Gupta2023} 


\begin{figure}[h!]
	\centering
	\includegraphics[width=0.4\textwidth]{Figuras/TIPO H.png}
	\caption{\textit{Geometría auxética tipo “H”.}
		Tomado de \textit{“High expansion auxetic skin graft simulants for severe burn injury mitigation”} 
		\parencite{Gupta2023}.}
	\label{fig:geometria_auxetica_h}
\end{figure}



\subsubsection{Geometría auxética tipo “RR” }
Estructura formada por rectángulos unidos en los vértices que tienen la capacidad de girar durante la deformación. Debido al anterior forma vacíos ovalados cuando se estira por lo cual, tiene una expansión moderada y efectiva sin pérdida estructural. Esta geometría presenta una relación de mallado de aproximadamente 2.8:1. \cite{Gupta2023} 

\begin{figure}[h!]
	\centering
	\includegraphics[width=0.4\textwidth]{Figuras/TIPO RR.png}
	\caption{\textit{Geometría auxética tipo “RR”.}
		Tomado de \textit{“High expansion auxetic skin graft simulants for severe burn injury mitigation”} 
		\parencite{Gupta2023}.}
	\label{fig:geometria_auxetica_rr}
\end{figure}



\subsubsection{Geometría auxética tipo “AS” }
Geometría formada por ranuras alternas cortadas en direcciones opuestas las cuales al expandirse forman vacíos circulares. Logra una expansión de área vacía de hasta 863\% y una relación de mallado de 1.5:1. Bajo altos esfuerzos puede comprometer la integridad de la estructura. \cite{Gupta2023} 


\begin{figure}[h!]
	\centering
	\includegraphics[width=0.4\textwidth]{Figuras/TIPO AS.png}
	\caption{\textit{Geometría auxética tipo “AS”.}
		Tomado de \textit{“High expansion auxetic skin graft simulants for severe burn injury mitigation”} 
		\parencite{Gupta2023}.}
	\label{fig:geometria_auxetica_as}
\end{figure}



 \subsection{Matriz de filtrado}
 
 
 \subsubsection{Matriz de filtrado del material}
 
%%%%%MATRIZ FILTRADO MATERIAL%%%%%%%%%%%%%%%%%%%%%%%%%%%%%%%%%%%%%%%%%%%%%%%%%%%%%%%%%%%%%%%
\newcolumntype{C}[1]{>{\centering\arraybackslash}m{#1}}
\setlength{\tabcolsep}{4pt}
\renewcommand{\arraystretch}{1.15}
\setlength\LTleft{\fill}
\setlength\LTright{\fill}

\begin{longtable}{C{5cm}cccccccc}
\caption{Matriz de filtrado de los conceptos del material}\\
\toprule
\textbf{Conceptos} &
\textbf{A} & \textbf{B} & \textbf{C} & \textbf{D} & \textbf{E} & \textbf{F} &
\textbf{Evaluación neta} & \textbf{¿Continúa?} \\
\midrule
\endfirsthead

\toprule
\textbf{Conceptos} &
\textbf{A} & \textbf{B} & \textbf{C} & \textbf{D} & \textbf{E} & \textbf{F} &
\textbf{Evaluación neta} & \textbf{¿Continúa?} \\
\midrule
\endhead

Referencia (Colágeno tipo I) & 0  & 0  & 0  & 0  & 0  & 0  & 0  & No \\
Piel de tilapia              & 1  & 1  & 0  & 0  & 0  & 0  & 2  & No \\
Quitosano                    & 0  & 0  & 1  & 1  & 1  & 1  & 4  & Sí \\
Alginato                     & -1 & 0  & 0  & 0  & -1 & 1  & -1 & Sí \\
Agarosa                      & -1 & 0  & 0  & 0  & -1 & 1  & -1 & Sí \\
Membrana amniótica           & 1  & 1  & 1  & 1  & 1  & -1 & 4  & Sí \\
PVP                          & -1 & 0  & 0  & 0  & -1 & 1  & -1 & No \\
Poliestireno                 & -1 & 1  & -1 & -1 & -1 & 1  & -1 & No \\
\bottomrule
\end{longtable}

\vspace{-0.4em}
\begin{center}
\begin{minipage}{0.96\textwidth} 
\centering
\textbf{Nota.} Para facilitar la lectura, los criterios de selección se codifican con letras (A–F) e indican entre paréntesis su peso relativo, cuyas ponderaciones suman el 100\%: \textbf{A}, mejora de la calidad de vida del paciente (13\%); \textbf{B}, mantenimiento de las propiedades durante el tiempo de aplicación (14\%); \textbf{C}, ausencia de citotoxicidad del material (25\%); \textbf{D}, promoción de la proliferación celular (16\%); \textbf{E}, contribución a una cicatrización adecuada (15\%); y \textbf{F}, facilidad de obtención–elaboración (17\%).

\end{minipage}
\end{center}




 \subsubsection{Matriz de filtrado del aditivo}
%%%%%%MATRIZ FILTRADO ADITIVO%%%%%%%%%%%%%%%%%%%%%%%%%%%%%%%%%%%%%%%%%%%%%%%%%%%%%%%%%%%%%%%%
% ===== MATRIZ FILTRADO ADITIVO =====
% (Mismo preámbulo que tu tabla anterior)
% \usepackage{booktabs}
% \usepackage{longtable}
% \usepackage{array}
\newcolumntype{C}[1]{>{\centering\arraybackslash}m{#1}}
\setlength{\tabcolsep}{4pt}
\renewcommand{\arraystretch}{1.15}
\setlength\LTleft{\fill}
\setlength\LTright{\fill}

\begin{longtable}{C{5cm}cccccc}
\caption{Matriz de filtrado de los conceptos de aditivo}\\
\toprule
\textbf{Conceptos} &
\textbf{A} & \textbf{B} & \textbf{C} & \textbf{D} &
\textbf{Evaluación neta} & \textbf{¿Continúa?} \\
\midrule
\endfirsthead

\toprule
\textbf{Conceptos} &
\textbf{A} & \textbf{B} & \textbf{C} & \textbf{D} &
\textbf{Evaluación neta} & \textbf{¿Continúa?} \\
\midrule
\endhead

Sulfadiazina de plata (Referencia) & 0  & 0  & 0  & 0  & 0 & No \\
Gentamicina                         & 1  & 0  & 0  & 1  & 2 & No \\
Minociclina                         & 1  & 0  & 0  & 0  & 1 & No \\
Vancomicina                         & 1  & 0  & 1  & 0  & 2 & No \\
Colagenasa                          & 1  & 1  & 1  & 0  & 3 & Sí \\
Fitostimoline                       & 1  & 1  & 1  & 1  & 4 & Sí \\
Ácido hialurónico                   & 1  & 1  & 0  & 1  & 3 & Sí \\
Aloe vera                           & -1 & -1 & 1  & 1  & 0 & No \\
Caléndula                           & -1 & -1 & 1  & 1  & 0 & No \\
\bottomrule
\end{longtable}

\vspace{-0.4em}
\begin{center}
\begin{minipage}{0.95\textwidth}
\centering
\textbf{Nota.} Para facilitar la lectura, los criterios de selección se codifican con letras (A–D) e incluyen entre paréntesis su peso relativo: 
\textbf{A}, mejora de la calidad de vida del paciente (19\%); 
\textbf{B}, biocompatibilidad del apósito (37\%); 
\textbf{C}, contribución a una buena cicatrización (21\%); 
\textbf{D}, facilidad de obtención–elaboración (23\%).
\end{minipage}
\end{center}

 \subsubsection{Matriz de filtrado de la geometría}

%%%%%%%MATRIZ FILTRADO DE GEOMETRIA%%%%%%%%%%%%%%%%%%%%%%%%%%%%%%%%%%%%%%%%%%%%%%%%%%%%%%%
% ===== MATRIZ FILTRADO GEOMETRÍA =====
% Usa el mismo preámbulo que ya tienes (booktabs, longtable, array y la col. C{})

\begin{longtable}{C{5.2cm}cccccc}
\caption{Matriz de filtrado de los conceptos de geometría}\\
\toprule
\textbf{Conceptos} &
\textbf{A} & \textbf{B} & \textbf{C} & \textbf{D} &
\textbf{Evaluación neta} & \textbf{¿Continúa?} \\
\midrule
\endfirsthead

\toprule
\textbf{Conceptos} &
\textbf{A} & \textbf{B} & \textbf{C} & \textbf{D} &
\textbf{Evaluación neta} & \textbf{¿Continúa?} \\
\midrule
\endhead

Referencia (Geometría tradicional de primer orden) & 0  & 0  & 0  & 0  & 0  & No \\
Tradicional de segundo orden                       & 1  & 0  & 0  & 0  & 1  & No \\
Tradicional de tercer orden                         & 1  & 0  & 0  & 0  & 1  & No \\
Auxética tipo ``Y''                                 & 1  & 0  & 1  & 0  & 2  & Sí \\
Auxética tipo ``I''                                 & 1  & 0  & 1  & 0  & 2  & Sí \\
Auxética tipo panal de abejas                       & 1  & 0  & 0  & 0  & 1  & No \\
Auxética tipo ``H''                                 & -1 & 0  & -1 & 0  & -2 & No \\
Auxética tipo ``RR''                                & 1  & 0  & 0  & 0  & 1  & No \\
Auxética tipo ``AS''                                & 0  & 0  & -1 & 0  & -1 & No \\
\bottomrule
\end{longtable}

\vspace{-0.4em}
\begin{center}
\begin{minipage}{0.95\textwidth}
\centering
\textbf{Nota.} Para facilitar la lectura, los criterios de selección se codifican con letras (A–D) y se indica entre paréntesis su peso relativo: \textbf{A}, mejora de la calidad de vida del paciente (19\,\%); \textbf{B}, biocompatibilidad del apósito (37\,\%); \textbf{C}, contribución a una buena cicatrización (21\,\%); \textbf{D}, facilidad de obtención–elaboración (23\,\%).

\end{minipage}
\end{center}



%%%%%%%%%%%%%%%%%%%%%%%%%%%%%%%%%%%%%%%%%%%%%%%%%%%%%%%%%%%%%%%%%%%%%%%%%%%%%
\subsubsection{Matriz de evaluación de concepto}

%%%%%%MATERIAL%%%%%

{\scriptsize
	\setlength{\tabcolsep}{2.3pt}   % menos espacio horizontal
	\renewcommand{\arraystretch}{1.1}
	\setlength\LTleft{\fill}
	\setlength\LTright{\fill}
	
	\begin{longtable}{C{3.0cm} *{12}{c} cc}
		\caption{Matriz de evaluación del concepto de material (con calificación y evaluación ponderada)}\\
		\toprule
		\multirow{3}{*}{\textbf{Conceptos}} &
		\multicolumn{12}{c}{\textbf{Criterios (A–F)}} &
		\multirow{3}{*}{\textbf{Eval. neta}} &
		\multirow{3}{*}{\textbf{¿Continúa?}} \\
		\cmidrule(lr){2-13}
		& \multicolumn{2}{c}{\textbf{A}} & \multicolumn{2}{c}{\textbf{B}} &
		\multicolumn{2}{c}{\textbf{C}} & \multicolumn{2}{c}{\textbf{D}} &
		\multicolumn{2}{c}{\textbf{E}} & \multicolumn{2}{c}{\textbf{F}} & & \\
		& \textit{Calif.} & \textit{Pond.} &
		\textit{Calif.} & \textit{Pond.} &
		\textit{Calif.} & \textit{Pond.} &
		\textit{Calif.} & \textit{Pond.} &
		\textit{Calif.} & \textit{Pond.} &
		\textit{Calif.} & \textit{Pond.} & & \\
		\midrule
		\endfirsthead
		
		\toprule
		\multirow{3}{*}{\textbf{Conceptos}} &
		\multicolumn{12}{c}{\textbf{Criterios (A–F)}} &
		\multirow{3}{*}{\textbf{Evaluación neta}} &
		\multirow{3}{*}{\textbf{¿Continúa?}} \\
		\cmidrule(lr){2-13}
		& \multicolumn{2}{c}{\textbf{A}} & \multicolumn{2}{c}{\textbf{B}} &
		\multicolumn{2}{c}{\textbf{C}} & \multicolumn{2}{c}{\textbf{D}} &
		\multicolumn{2}{c}{\textbf{E}} & \multicolumn{2}{c}{\textbf{F}} & & \\
		& \textit{Calif.} & \textit{Pond.} &
		\textit{Calif.} & \textit{Pond.} &
		\textit{Calif.} & \textit{Pond.} &
		\textit{Calif.} & \textit{Pond.} &
		\textit{Calif.} & \textit{Pond.} &
		\textit{Calif.} & \textit{Pond.} & & \\
		\midrule
		\endhead
		
		Referencia (Colágeno tipo I) & 3 & 0.39 & 3 & 0.42 & 3 & 0.75 & 3 & 0.48 & 3 & 0.45 & 3 & 0.51 & 3.00 & No \\
		Piel de tilapia              & 4 & 0.52 & 4 & 0.56 & 3 & 0.75 & 3 & 0.48 & 3 & 0.45 & 3 & 0.51 & 3.27 & No \\
		Quitosano                    & 3 & 0.39 & 3 & 0.42 & 4 & 1.00 & 4 & 0.64 & 4 & 0.60 & 4 & 0.68 & 3.73 & Sí \\
		Alginato                     & 3 & 0.39 & 3 & 0.42 & 3 & 0.75 & 3 & 0.48 & 2 & 0.30 & 4 & 0.68 & 3.02 & Sí \\
		Agarosa                      & 3 & 0.39 & 3 & 0.42 & 3 & 0.75 & 3 & 0.48 & 2 & 0.30 & 4 & 0.68 & 3.02 & Sí \\
		Membrana amniótica           & 4 & 0.52 & 4 & 0.56 & 4 & 1.00 & 4 & 0.64 & 4 & 0.60 & 1 & 0.17 & 3.49 & Sí \\
		PVP                          & 2 & 0.26 & 3 & 0.42 & 3 & 0.75 & 3 & 0.48 & 2 & 0.30 & 4 & 0.68 & 2.89 & No \\
		Poliestireno                 & 2 & 0.26 & 4 & 0.56 & 2 & 0.50 & 2 & 0.32 & 2 & 0.30 & 4 & 0.68 & 2.62 & No \\
		\bottomrule
	\end{longtable}
} % fin \scriptsize

\vspace{-0.4em}
\begin{center}
	\begin{minipage}{0.95\textwidth}
		\centering
		\normalsize
		\textbf{Nota.} Para facilitar la lectura, los criterios se codifican con letras (A–F) y se indica entre paréntesis su peso relativo; las columnas dobles muestran la \textit{calificación} (1–4) y la \textit{evaluación ponderada} (calificación $\times$ peso): 
		\textbf{A}, mejora de la calidad de vida del paciente (13\,\%); 
		\textbf{B}, mantenimiento de propiedades durante el tiempo de aplicación (14\,\%); 
		\textbf{C}, ausencia de citotoxicidad (25\,\%); 
		\textbf{D}, promoción de la proliferación celular (16\,\%); 
		\textbf{E}, contribución a una buena cicatrización (15\,\%); 
		\textbf{F}, facilidad de obtención–elaboración (17\,\%). En cada criterio, \textit{Calif.} corresponde a la calificación (escala 1–5) y \textit{Pond.} al producto \emph{calificación $\times$ peso}; la \textbf{Evaluación neta} es la suma de todas las ponderaciones. 
		La escala de calificación se define así: \textbf{1} = mucho peor que la referencia; \textbf{2} = peor que la referencia; \textbf{3} = igual que la referencia; \textbf{4} = mejor que la referencia; \textbf{5} = mucho mejor que la referencia.
	\end{minipage}
\end{center}










%%%%%ADITIVO%%%%%

%%%%% ADITIVO %%%%%
\sisetup{locale=DE, table-number-alignment=center, table-text-alignment=center}

\newcolumntype{C}[1]{>{\centering\arraybackslash}m{#1}}
\newcommand{\hdr}[1]{\multicolumn{1}{c}{#1}}

{\scriptsize
	\setlength{\tabcolsep}{2.3pt}
	\renewcommand{\arraystretch}{1.1}
	\setlength\LTleft{\fill}
	\setlength\LTright{\fill}
	
	\begin{longtable}{
			C{3.2cm}
			S[table-format=1.0]  S[table-format=1.2]
			S[table-format=1.0]  S[table-format=1.2]
			S[table-format=1.0]  S[table-format=1.0]  S[table-format=1.0]  S[table-format=1.2]
			S[table-format=1.0]  S[table-format=1.2]
			S[table-format=1.2]  c
		}
		\caption{Matriz de evaluación del concepto \textit{aditivo} (con calificación y evaluación ponderada)}\\
		\toprule
		\multicolumn{1}{c}{} &
		\multicolumn{11}{c}{\textbf{Criterios}} &
		\multicolumn{1}{c}{}\\
		\cmidrule(lr){2-12}
		\textbf{Conceptos} &
		\multicolumn{2}{c}{\makecell{\textbf{A}\\(13\%)}} &
		\multicolumn{2}{c}{\makecell{\textbf{B}\\(20\%)}} &
		\multicolumn{4}{c}{\textbf{C}} &
		\multicolumn{2}{c}{\makecell{\textbf{D}\\(15\%)}} &
		\textbf{Eval. neta} & \textbf{¿Continúa?} \\
		\cmidrule(lr){2-3}\cmidrule(lr){4-5}\cmidrule(lr){6-9}\cmidrule(lr){10-11}
		& \hdr{\textit{Calif.}} & \hdr{\textit{Pond.}} &
		\hdr{\textit{Calif.}} & \hdr{\textit{Pond.}} &
		\hdr{\makecell{\textit{1.\,Calif.}\\(25\%)}} &
		\hdr{\makecell{\textit{2.\,Calif.}\\(10\%)}} &
		\hdr{\makecell{\textit{3.\,Calif.}\\(17\%)}} &
		\hdr{\textit{Pond.}} &
		\hdr{\textit{Calif.}} & \hdr{\textit{Pond.}} &
		\hdr{} & \hdr{} \\
		\midrule
		\multicolumn{1}{c}{\textbf{Peso (\%)}} &
		\multicolumn{2}{c}{13} &
		\multicolumn{2}{c}{20} &
		\multicolumn{1}{c}{25} & \multicolumn{1}{c}{10} & \multicolumn{1}{c}{17} & \multicolumn{1}{c}{n/a} &
		\multicolumn{2}{c}{15} &
		\multicolumn{2}{c}{} \\
		\midrule
		\endfirsthead
		
		\toprule
		\multicolumn{1}{c}{} &
		\multicolumn{11}{c}{\textbf{Criterios}} &
		\multicolumn{1}{c}{}\\
		\cmidrule(lr){2-12}
		\textbf{Conceptos} &
		\multicolumn{2}{c}{\makecell{\textbf{A}\\(13\%)}} &
		\multicolumn{2}{c}{\makecell{\textbf{B}\\(20\%)}} &
		\multicolumn{4}{c}{\textbf{C}} &
		\multicolumn{2}{c}{\makecell{\textbf{D}\\(15\%)}} &
		\textbf{Eval. neta} & \textbf{¿Continúa?} \\
		\cmidrule(lr){2-3}\cmidrule(lr){4-5}\cmidrule(lr){6-9}\cmidrule(lr){10-11}
		& \hdr{\textit{Calif.}} & \hdr{\textit{Pond.}} &
		\hdr{\textit{Calif.}} & \hdr{\textit{Pond.}} &
		\hdr{\makecell{\textit{1.\,Calif.}\\(25\%)}} &
		\hdr{\makecell{\textit{2.\,Calif.}\\(10\%)}} &
		\hdr{\makecell{\textit{3.\,Calif.}\\(17\%)}} &
		\hdr{\textit{Pond.}} &
		\hdr{\textit{Calif.}} & \hdr{\textit{Pond.}} &
		\hdr{} & \hdr{} \\
		\midrule
		\multicolumn{1}{c}{\textbf{Peso (\%)}} &
		\multicolumn{2}{c}{13} &
		\multicolumn{2}{c}{20} &
		\multicolumn{1}{c}{25} & \multicolumn{1}{c}{10} & \multicolumn{1}{c}{17} & \multicolumn{1}{c}{n/a} &
		\multicolumn{2}{c}{15} &
		\multicolumn{2}{c}{} \\
		\midrule
		\endhead
		
		% ===== DATOS =====
		Sulfadiazina de plata (Referencia) & 3 & 0,39 & 3 & 0,60 & 3 & 3 & 3 & 1,56 & 3 & 0,45 & 4,44 & No \\
		Gentamicina                        & 3 & 0,39 & 3 & 0,60 & 2 & 3 & 2 & 1,14 & 3 & 0,45 & 4,44 & No \\
		Minociclina                        & 3 & 0,39 & 3 & 0,60 & 2 & 3 & 3 & 1,31 & 2 & 0,30 & 4,29 & No \\
		Vancomicina                        & 3 & 0,39 & 3 & 0,60 & 4 & 3 & 3 & 1,81 & 3 & 0,45 & 4,44 & No \\
		Colagenasa                         & 4 & 0,52 & 3 & 0,60 & 5 & 5 & 5 & 2,60 & 2 & 0,30 & 6,42 & Sí \\
		Fitostimoline                      & 4 & 0,52 & 3 & 0,60 & 5 & 5 & 5 & 2,60 & 3 & 0,45 & 6,57 & Sí \\
		Ácido hialurónico                  & 4 & 0,52 & 3 & 0,60 & 5 & 5 & 5 & 2,60 & 3 & 0,45 & 6,57 & Sí \\
		Aloe vera                          & 2 & 0,26 & 2 & 0,40 & 2 & 1 & 1 & 0,77 & 4 & 0,60 & 2,26 & No \\
		Caléndula                          & 2 & 0,26 & 2 & 0,40 & 2 & 2 & 1 & 0,87 & 4 & 0,60 & 3,26 & No \\
		\bottomrule
	\end{longtable}
} % fin scriptsize

\begin{center} 
\begin{minipage}{0.96\textwidth} \centering \textbf{Nota.} Para facilitar la lectura, los criterios de selección se codifican con letras (A–D) e indican entre paréntesis su peso relativo: \textbf{A} (13\,\%): mejora de la calidad de vida del paciente; \textbf{B} (20\,\%): biocompatibilidad del apósito; \textbf{C}: Ayuda a una buena cicatrización, donde dicho criterio esta compuesto en tres subcriterios numerados —\textbf{1} (25\,\%): el apósito induce la \textit{granulación} de la herida; \textbf{2} (10\,\%): el apósito se adapta correctamente a \textit{pliegues y articulaciones}; \textbf{3} (17\,\%): el apósito presenta \textit{liberación controlada de fármacos}— cuya contribución se integra en la columna \textit{Pond.}; \textbf{D} (15\,\%): facilidad de obtención–elaboración. En cada criterio, \textit{Calif.} corresponde a la calificación (escala 1–5) y \textit{Pond.} al producto \emph{calificación $\times$ peso}; la \textbf{Evaluación neta} es la suma de todas las ponderaciones. La escala de calificación se define así: \textbf{1} = mucho peor que la referencia; \textbf{2} = peor que la referencia; \textbf{3} = igual que la referencia; \textbf{4} = mejor que la referencia; \textbf{5} = mucho mejor que la referencia. 
\end{minipage} 
\end{center}


%%%%%%%%%%%%%%%%%%%%%%%%%%%%%%%%%%%%%%%%%%%%%GEOMETRIA %%%%%%%%%%%%%%%%%%%%%%%%%%%%%

\sisetup{
	output-decimal-marker = {,},
	table-number-alignment = center,
	table-text-alignment   = center
}

% Columna centrada de ancho fijo (para "Conceptos")
\newcolumntype{C}[1]{>{\centering\arraybackslash}m{#1}}

% Definir hdr solo si no existe
\providecommand{\hdr}[1]{\multicolumn{1}{c}{#1}}

{\scriptsize
	\setlength{\tabcolsep}{2.3pt}
	\renewcommand{\arraystretch}{1.1}
	\setlength\LTright{\fill}
	
	\begin{longtable}{
			C{3.6cm}
			S[table-format=1.0]  S[table-format=1.2]
			S[table-format=1.0]  S[table-format=1.2]
			S[table-format=1.0]  S[table-format=1.0]  S[table-format=1.0]  S[table-format=1.2]
			S[table-format=1.0]  S[table-format=1.2]
			S[table-format=1.2]  c
		}
		\caption{Matriz de evaluación del concepto \textit{geometría} (con calificación y evaluación ponderada)}\\
		\toprule
		\multicolumn{1}{c}{} &
		\multicolumn{11}{c}{\textbf{Criterios}} &
		\multicolumn{1}{c}{}\\
		\cmidrule(lr){2-12}
		\textbf{Conceptos} &
		\multicolumn{2}{c}{\makecell{\textbf{A}\\(19\%)}} &
		\multicolumn{2}{c}{\makecell{\textbf{B}\\(37\%)}} &
		\multicolumn{4}{c}{\textbf{C}} &
		\multicolumn{2}{c}{\makecell{\textbf{D}\\(23\%)}} &
		\textbf{Eval. neta} & \textbf{¿Continúa?} \\
		\cmidrule(lr){2-3}\cmidrule(lr){4-5}\cmidrule(lr){6-9}\cmidrule(lr){10-11}
		& \hdr{\textit{Calif.}} & \hdr{\textit{Pond.}} &
		\hdr{\textit{Calif.}} & \hdr{\textit{Pond.}} &
		\hdr{\makecell{\textit{1.\,Calif.}\\(25\%)}} &
		\hdr{\makecell{\textit{2.\,Calif.}\\(10\%)}} &
		\hdr{\makecell{\textit{3.\,Calif.}\\(17\%)}} &
		\hdr{\textit{Pond.}} &
		\hdr{\textit{Calif.}} & \hdr{\textit{Pond.}} &
		\hdr{} & \hdr{} \\
		\midrule
		\multicolumn{1}{c}{\textbf{Peso (\%)}} &
		\multicolumn{2}{c}{19} &
		\multicolumn{2}{c}{37} &
		\multicolumn{1}{c}{25} & \multicolumn{1}{c}{10} & \multicolumn{1}{c}{17} & \multicolumn{1}{c}{n/a} &
		\multicolumn{2}{c}{23} &
		\multicolumn{2}{c}{} \\
		\midrule
		\endfirsthead
		
		\toprule
		\multicolumn{1}{c}{} &
		\multicolumn{11}{c}{\textbf{Criterios}} &
		\multicolumn{1}{c}{}\\
		\cmidrule(lr){2-12}
		\textbf{Conceptos} &
		\multicolumn{2}{c}{\makecell{\textbf{A}\\(19\%)}} &
		\multicolumn{2}{c}{\makecell{\textbf{B}\\(37\%)}} &
		\multicolumn{4}{c}{\textbf{C}} &
		\multicolumn{2}{c}{\makecell{\textbf{D}\\(23\%)}} &
		\textbf{Eval. neta} & \textbf{¿Continúa?} \\
		\cmidrule(lr){2-3}\cmidrule(lr){4-5}\cmidrule(lr){6-9}\cmidrule(lr){10-11}
		& \hdr{\textit{Calif.}} & \hdr{\textit{Pond.}} &
		\hdr{\textit{Calif.}} & \hdr{\textit{Pond.}} &
		\hdr{\makecell{\textit{1.\,Calif.}\\(25\%)}} &
		\hdr{\makecell{\textit{2.\,Calif.}\\(10\%)}} &
		\hdr{\makecell{\textit{3.\,Calif.}\\(17\%)}} &
		\hdr{\textit{Pond.}} &
		\hdr{\textit{Calif.}} & \hdr{\textit{Pond.}} &
		\hdr{} & \hdr{} \\
		\midrule
		\multicolumn{1}{c}{\textbf{Peso (\%)}} &
		\multicolumn{2}{c}{19} &
		\multicolumn{2}{c}{37} &
		\multicolumn{1}{c}{25} & \multicolumn{1}{c}{10} & \multicolumn{1}{c}{17} & \multicolumn{1}{c}{n/a} &
		\multicolumn{2}{c}{23} &
		\multicolumn{2}{c}{} \\
		\midrule
		\endhead
		
		% === Datos ===
		Referencia (Geometría tradicional de primer orden) & 3 & 0,57 & 3 & 1,11 & 3 & 3 & 3 & 1,56 & 3 & 0,69 & 3,93 & No \\
		Tradicional de segundo orden                       & 3 & 0,57 & 3 & 1,11 & 3 & 3 & 3 & 1,56 & 3 & 0,69 & 3,93 & No \\
		Tradicional de tercer orden                        & 3 & 0,57 & 3 & 1,11 & 3 & 4 & 3 & 1,81 & 3 & 0,69 & 4,18 & No \\
		Auxética tipo Y                                    & 4 & 0,76 & 3 & 1,11 & 3 & 4 & 3 & 1,81 & 4 & 0,92 & 4,60 & No \\
		Auxética tipo I                                    & 4 & 0,76 & 3 & 1,11 & 4 & 4 & 3 & 1,98 & 4 & 0,92 & 4,77 & Sí \\
		Auxética tipo panal de abejas                      & 4 & 0,76 & 3 & 1,11 & 3 & 3 & 3 & 1,56 & 3 & 0,69 & 4,12 & No \\
		Auxética tipo H                                    & 2 & 0,38 & 3 & 1,11 & 3 & 2 & 3 & 1,31 & 3 & 0,69 & 3,49 & No \\
		Auxética tipo RR                                   & 4 & 0,76 & 3 & 1,11 & 2 & 2 & 3 & 1,14 & 2 & 0,46 & 3,47 & No \\
		Auxética tipo AS                                   & 3 & 0,57 & 3 & 1,11 & 2 & 3 & 3 & 1,14 & 3 & 0,69 & 3,51 & No \\
		\bottomrule
	\end{longtable}
} % fin scriptsize

\begin{center}
	\begin{minipage}{0.96\textwidth}
		\centering
		\textbf{Nota.} Para facilitar la lectura, los criterios de selección se codifican con letras (A–D)… 
	\end{minipage}
\end{center}


\begin{center}
\begin{minipage}{0.96\textwidth}
\centering
\textbf{Nota.} Para facilitar la lectura, los criterios de selección se codifican con letras (A–D) e indican entre paréntesis su peso relativo: 
\textbf{A} (19\,\%): mejora de la calidad de vida del paciente; 
\textbf{B} (37\,\%): biocompatibilidad del apósito; 
\textbf{C}: criterio compuesto con tres subcriterios —\textbf{1} (25\,\%): inducción de \textit{granulación} de la herida; \textbf{2} (10\,\%): adaptación a \textit{pliegues y articulaciones}; \textbf{3} (17\,\%): \textit{liberación controlada de fármacos}— cuya contribución se integra en la columna \textit{Pond.}; 
\textbf{D} (23\,\%): facilidad de obtención–elaboración. 
En cada criterio, \textit{Calif.} es la calificación (escala 1–5) y \textit{Pond.} el producto \emph{calificación $\times$ peso}; la \textbf{Evaluación neta} es la suma de todas las ponderaciones. 
La escala utilizada es: \textbf{1} = mucho peor que la referencia; \textbf{2} = peor que la referencia; \textbf{3} = igual que la referencia; \textbf{4} = mejor que la referencia; \textbf{5} = mucho mejor que la referencia.
\end{minipage}
\end{center}


  \section{Selección del Concepto}
  Teniendo en cuenta la evaluación neta obtenida en las matrices para cada aspecto: material, aditivo y geometría se seleccionaron aquellos conceptos que obtuvieron el mayor valor. Por ende, se plantea el desarrollo de un apósito diseñado con  geometría auxética tipo “I” , esta estructura se caracteriza por presentar un coeficiente de poisson de hasta -1.18 lo cual indica una alta capacidad de expansión. Además, a partir de la revisión de literatura se evidencia que esta geometría tiene una alta relación de mallado comparado con las otras geometrías propuestas en los conceptos, esto indica que la estructura es ideal para mitigar la limitación actual de las geometrías tradicionales de apósitos. Es importante resaltar que se deben evaluar diferentes alternativas de la geometría seleccionada con el fin de obtener la mejor distribución del patrón, es decir que, se propone realizar pruebas de simulación donde se diseñe la geometría auxética tipo “I” variando la longitud de lado y la separación entre cada patrón estructural. 
  
  Además, las propiedades mecánicas y funcionales del apósito están estrechamente ligadas al material seleccionado. En este sentido, tras el análisis comparativo de los conceptos propuestos, el quitosano fue el material que obtuvo la mayor puntuación, destacándose por sus sobresalientes propiedades fisicoquímicas y biológicas. El quitosano es un polisacárido natural derivado de la quitina, conocido por su biocompatibilidad, biodegradabilidad, propiedades antimicrobianas y su capacidad para favorecer la proliferación celular y la regeneración tisular. Sin embargo, reconociendo las limitaciones prácticas del proyecto particularmente en lo referente a la disponibilidad y facilidad de obtención de los materiales, se propone la evaluación de materiales híbridos en las que el quitosano se combine con alginato y agarosa, con el propósito de potenciar sus propiedades y superar posibles deficiencias individuales. Estas combinaciones permitirán explorar mejoras en términos de resistencia mecánica, capacidad de absorción, estabilidad estructural y comportamiento biológico, contribuyendo así a la optimización integral del apósito propuesto.
  
  Por último, tras el análisis comparativo de los conceptos propuestos para el aditivo mediante una matriz de filtrado, se identificó que fitostimuline, ácido hialurónico y colagenasa obtuvieron las mejores puntuaciones, destacándose por sus propiedades biológicas, capacidad para favorecer la cicatrización y su disponibilidad. Estos aditivos demostraron un alto potencial en la estimulación de la granulación y epitelización, biocompatibilidad y adecuada adaptación a zonas de pliegues y articulaciones. Sin embargo, la colagenasa presenta una desventaja en cuanto a su disponibilidad, ya que actualmente solo se comercializa en Estados Unidos, lo que limita su accesibilidad y viabilidad dentro del contexto del proyecto. Por ello, se propone priorizar la evaluación de fitostimoline y ácido hialurónico, considerando su fácil obtención en el mercado local y su capacidad para combinarse con el material del apósito, lo que permitiría optimizar las propiedades terapéuticas y funcionales del apósito propuesto, garantizando una solución eficaz, accesible y de alto desempeño biológico.
 
 
 
    \subsubsection*{Concepto de material}

    \subsubsection{Concepto de aditivo}
    
    \subsubsection{Concepto de geometría}
    
    A partir del concepto seleccionado para el aspecto de geometría, se propuso la simulación de la geometría auxética tipo “I”, variando la longitud de los bordes y la separación entre las repeticiones del patrón. De este modo, se planteó el desarrollo de la simulación mediante un análisis de elementos finitos, con el fin de obtener resultados sobre la deformación unitaria, la tensión y el desplazamiento, y así definir el diseño con mejor cubrimiento de área, basado en el comportamiento auxético de cada estructura para el desarrollo del apósito.
    
    
\begin{figure}[h!]
	\centering
	\includegraphics[width=0.4\textwidth]{Figuras/MODIFICACIONESATIPOI.png}
	\caption{\textit{Parámetros variados en la geometría auxética tipo “I”.}
		Imagen modificada a partir de \textit{“High expansion auxetic skin graft simulants for severe burn injury mitigation”} 
		\parencite{Gupta2023}.}
	\label{fig:geometria_auxetica_i_modificada}
\end{figure}





\section{Discusión}


\end{document}
